\documentclass[11pt,a4paper]{article}

% ── Codificación, fuentes y lenguaje ───────────────────────────────────────────
\usepackage[utf8]{inputenc}
\usepackage[T1]{fontenc}
\usepackage[default,scale=0.95]{sourcesanspro}
\renewcommand{\familydefault}{\sfdefault}
\usepackage[spanish,es-noshorthands,es-tabla]{babel}

% ── Geometría y espaciado ─────────────────────────────────────────────────────
\usepackage[top=2.2cm, bottom=2.5cm, left=2.2cm, right=2.2cm, headheight=15pt]{geometry}
\usepackage{setspace}
\setstretch{1.18}
\usepackage{parskip}

% ── Matemáticas y unidades ────────────────────────────────────────────────────
\usepackage{amsmath}
\usepackage{amssymb}
\usepackage{siunitx}

% ── Circuitos ─────────────────────────────────────────────────────────────────
\usepackage[european, straightvoltages]{circuitikz}

% ── Tablas ────────────────────────────────────────────────────────────────────
\usepackage{booktabs}
\usepackage{tabularx}
\usepackage{array}
\usepackage{longtable}
\usepackage{colortbl}
\usepackage{enumitem}

% ── Colores, cajas e iconos ───────────────────────────────────────────────────
\usepackage[dvipsnames]{xcolor}
\usepackage{tcolorbox}
\tcbuselibrary{skins, breakable}
\usepackage{fontawesome5}
\usepackage{graphicx}

% ── Listados de código (dark-mode infográfico) ────────────────────────────────
%   Estilo base: fondo oscuro con números de línea. Los bloques lstlisting se
%   envuelven automáticamente en una caja tcolorbox redondeada.
\usepackage{listings}

% ── Secciones con barra de color (infografía) ────────────────────────────────
\usepackage{titlesec}
\titleformat{\section}
  {\normalfont\LARGE\bfseries\color{azulNoche}}
  {\colorbox{cyanNeon}{\color{fondoOscuro}\thesection}}{0.7em}{}
  [\vspace{2pt}\color{cyanNeon}\rule{\linewidth}{1.8pt}]
\titlespacing*{\section}{0pt}{24pt}{10pt}
\titleformat{\subsection}
  {\normalfont\large\bfseries\color{azulNoche}}
  {\textcolor{cyanNeon}{\thesubsection}}{0.7em}{}
  [\vspace{1pt}\color{grisLinea}\rule{\linewidth}{0.6pt}]
\titlespacing*{\subsection}{0pt}{14pt}{6pt}
\titleformat{\subsubsection}
  {\normalfont\normalsize\bfseries\color{azulNoche}}
  {\textcolor{cyanNeon}{\thesubsubsection}}{0.7em}{}
\titlespacing*{\subsubsection}{0pt}{10pt}{4pt}

% ── Cabeceras y pies de página ────────────────────────────────────────────────
\usepackage{fancyhdr}
\usepackage{lastpage}
\pagestyle{fancy}
\fancyhf{}
\renewcommand{\headrulewidth}{0pt}
\renewcommand{\footrulewidth}{0pt}
\fancyfoot[C]{%
  \begin{minipage}{\linewidth}
    \vspace{2pt}
    \color{cyanNeon}\rule{\linewidth}{1.2pt}\\[2pt]
    \centering\color{grisTexto}\footnotesize
    Página \thepage\ de \pageref{LastPage}
  \end{minipage}%
}

% ── Hiperenlaces ──────────────────────────────────────────────────────────────
\usepackage[colorlinks=true, linkcolor=azulNoche, urlcolor=cyanNeon]{hyperref}
\sloppy
\emergencystretch 3em

% ── Paleta de colores — tecnológico dark-mode ──────────────────────────────────
\definecolor{fondoOscuro}{HTML}{0D1B2A}%             Dark navy — fondo banda título
\definecolor{verdeTurquesa}{HTML}{004D40}%           Verde turquesa — bandas de marca
\definecolor{azulNoche}{HTML}{1B3A6B}%               Midnight blue — secciones, marcos
\definecolor{azulMedio}{HTML}{2A5298}%               Azul medio — variante de acento
\definecolor{cyanNeon}{HTML}{00C8E0}%                Cyan eléctrico — acentos primarios
\definecolor{verdeSignal}{HTML}{00B887}%             Verde señal — fortalezas, éxito
\definecolor{naranjaVivo}{HTML}{FF7043}%             Naranja vivo — mejoras, advertencias
\definecolor{rojoAlerta}{HTML}{E53935}%              Rojo alerta — errores
\definecolor{grisPapel}{HTML}{F4F6FA}%               Fondo tarjetas claras
\definecolor{grisLinea}{HTML}{C8D0E0}%               Bordes sutiles
\definecolor{grisTexto}{HTML}{6B7A99}%               Texto secundario
\definecolor{amarilloNota}{HTML}{FFB300}%            Notas / advertencias doradas
\definecolor{codigoFondo}{HTML}{1E2D3D}%             Fondo del bloque de código
\definecolor{codigoComentario}{HTML}{7A8FA6}%        Comentarios
\definecolor{codigoCadena}{HTML}{FF9D5C}%            Cadenas / strings
\definecolor{codigoKeyword}{HTML}{00C8E0}%           Palabras clave
\definecolor{codigoNumero}{HTML}{00E5A0}%            Números

% ── Banda de título: portada infográfica ──────────────────────────────────────
%   Diseño en 2 capas: banda de color (por defecto fondo oscuro) + franja de
%   acento cyan abajo. Uso: \bandaTitulo[color]{título}{subtítulo}.
%   El icono se puede cambiar con \renewcommand{\iconoBanda}{...} (FontAwesome).
\newcommand{\iconoBanda}{microchip}
\newcommand{\bandaTitulo}[3][fondoOscuro]{%
  \begin{tcolorbox}[
    enhanced,
    colback=#1, colframe=#1,
    boxrule=0pt, arc=8pt, outer arc=8pt,
    width=\linewidth,
    top=18pt, bottom=6pt, left=14pt, right=14pt,
    borderline south={3.5pt}{0pt}{cyanNeon},
  ]
    \begin{minipage}[c]{0.07\linewidth}
      \centering
      \textcolor{cyanNeon}{\fontsize{30}{30}\selectfont\faIcon{\iconoBanda}}
    \end{minipage}%
    \hspace{8pt}%
    \begin{minipage}[c]{0.88\linewidth}
      {\fontsize{20}{24}\selectfont\bfseries\color{white}#2}\par\vspace{5pt}%
      \textcolor{cyanNeon!80!white}{\small\faIcon{angle-right}~#3}%
    \end{minipage}
    \vspace{4pt}
  \end{tcolorbox}%
  \vspace{8pt}%
}

% ── Tarjetas de datos (infografía) ────────────────────────────────────────────
\tcbset{
  tarjetaDato/.style={
    enhanced, breakable,
    arc=6pt, outer arc=6pt,
    colback=grisPapel, colframe=azulNoche,
    boxrule=1pt,
    fonttitle=\bfseries\small, coltitle=white,
    attach boxed title to top left={yshift=-3mm, xshift=6mm},
    boxed title style={
      colback=azulNoche, colframe=azulNoche,
      arc=4pt, boxrule=0pt,
    },
    drop fuzzy shadow=azulNoche!30!white,
    top=8pt, bottom=8pt, left=10pt, right=10pt,
  },
  tarjetaIcono/.style={
    enhanced, breakable,
    arc=6pt, outer arc=6pt,
    colback=white, colframe=grisLinea, boxrule=0.8pt,
    fonttitle=\bfseries\small, coltitle=azulNoche,
    borderline west={3pt}{0pt}{cyanNeon},
    drop fuzzy shadow=azulNoche!25!white,
    top=6pt, bottom=6pt, left=10pt, right=10pt,
  },
  cajaTitulo/.style={
    enhanced, breakable,
    arc=6pt, outer arc=6pt,
    colback=azulNoche!10!grisPapel, colframe=azulNoche,
    boxrule=1pt,
    fonttitle=\bfseries, coltitle=white,
    attach boxed title to top left={yshift=-3mm, xshift=6mm},
    boxed title style={
      colback=azulNoche, colframe=azulNoche,
      arc=4pt, boxrule=0pt,
    },
    drop fuzzy shadow=azulNoche!25!white,
    top=8pt, bottom=8pt, left=10pt, right=10pt,
  },
  cajaContenido/.style={
    enhanced, breakable,
    arc=5pt, outer arc=5pt,
    colback=grisPapel, colframe=grisLinea,
    boxrule=0.8pt,
    fonttitle=\bfseries\small, coltitle=azulNoche,
    top=6pt, bottom=6pt, left=10pt, right=10pt,
  },
  cajaFortaleza/.style={
    enhanced, breakable,
    arc=6pt, outer arc=6pt,
    colback=verdeSignal!8!white, colframe=verdeSignal,
    boxrule=1pt,
    borderline west={4pt}{0pt}{verdeSignal},
    fonttitle=\bfseries\small, coltitle=verdeSignal!70!black,
    drop fuzzy shadow=verdeSignal!20!white,
    top=6pt, bottom=6pt, left=10pt, right=10pt,
  },
  cajaMejora/.style={
    enhanced, breakable,
    arc=6pt, outer arc=6pt,
    colback=naranjaVivo!8!white, colframe=naranjaVivo,
    boxrule=1pt,
    borderline west={4pt}{0pt}{naranjaVivo},
    fonttitle=\bfseries\small, coltitle=naranjaVivo!80!black,
    drop fuzzy shadow=naranjaVivo!20!white,
    top=6pt, bottom=6pt, left=10pt, right=10pt,
  },
  cajaRecomendacion/.style={
    enhanced, breakable,
    arc=6pt, outer arc=6pt,
    colback=cyanNeon!6!white, colframe=cyanNeon!60!azulNoche,
    boxrule=1pt,
    borderline west={4pt}{0pt}{cyanNeon},
    fonttitle=\bfseries\small, coltitle=azulNoche,
    drop fuzzy shadow=azulNoche!20!white,
    top=6pt, bottom=6pt, left=10pt, right=10pt,
  },
}

% ── Ícono + texto (encabezados de sección y elementos) ────────────────────────
\newcommand{\iconotexto}[2]{\textcolor{cyanNeon}{\faIcon{#1}}~#2}

% ── Listados de código: estilo dark + auto-envoltura ─────────────────────────
\lstdefinestyle{estiloCodigo}{
    backgroundcolor=\color{codigoFondo},
    basicstyle=\ttfamily\small\color{white},
    commentstyle=\color{codigoComentario}\itshape,
    stringstyle=\color{codigoCadena},
    keywordstyle=\color{codigoKeyword}\bfseries,
    numberstyle=\tiny\color{codigoComentario},
    numbers=left,
    numbersep=10pt,
    showstringspaces=false,
    breaklines=true,
    breakatwhitespace=false,
    tabsize=4,
    frame=none,
    captionpos=b,
    aboveskip=6pt,
    belowskip=4pt,
    xleftmargin=14pt,
    framexleftmargin=14pt,
    literate=
      {á}{{\'a}}1 {é}{{\'e}}1 {í}{{\'i}}1 {ó}{{\'o}}1 {ú}{{\'u}}1
      {Á}{{\'A}}1 {É}{{\'E}}1 {Í}{{\'I}}1 {Ó}{{\'O}}1 {Ú}{{\'U}}1
      {ñ}{{\~n}}1 {Ñ}{{\~N}}1 {ü}{{\"u}}1 {Ü}{{\"U}}1
      {¿}{{?`}}1 {¡}{{!`}}1
}
\lstdefinestyle{estiloPython}{
    style=estiloCodigo,
    language=Python,
}
\lstset{style=estiloCodigo}
\BeforeBeginEnvironment{lstlisting}{%
  \begin{tcolorbox}[
    enhanced, breakable,
    arc=6pt, outer arc=6pt,
    colback=codigoFondo, colframe=azulNoche,
    boxrule=1pt,
    drop fuzzy shadow=azulNoche!25!white,
    top=2pt, bottom=2pt, left=0pt, right=0pt,
  ]%
}
\AfterEndEnvironment{lstlisting}{\end{tcolorbox}}

% ── Cajas de contenido temáticas ──────────────────────────────────────────────
\newtcolorbox{cajaRecuerda}{
    enhanced, breakable,
    arc=6pt, outer arc=6pt,
    colback=azulNoche!10!grisPapel,
    colframe=azulNoche, boxrule=1pt,
    borderline west={4pt}{0pt}{cyanNeon},
    fonttitle=\bfseries\small\color{white},
    title={\faIcon{info-circle}~Recuerda},
    attach boxed title to top left={yshift=-3mm, xshift=6mm},
    boxed title style={colback=azulNoche, arc=4pt, boxrule=0pt},
    drop fuzzy shadow=azulNoche!25!white,
    left=10pt, right=10pt, top=8pt, bottom=8pt,
}

\newtcolorbox{cajaConcepto}{
    enhanced, breakable,
    arc=6pt, outer arc=6pt,
    colback=verdeSignal!8!white,
    colframe=verdeSignal, boxrule=1pt,
    borderline west={4pt}{0pt}{verdeSignal},
    fonttitle=\bfseries\small\color{white},
    title={\faIcon{lightbulb}~Concepto clave},
    attach boxed title to top left={yshift=-3mm, xshift=6mm},
    boxed title style={colback=verdeSignal!80!black, arc=4pt, boxrule=0pt},
    drop fuzzy shadow=verdeSignal!20!white,
    left=10pt, right=10pt, top=8pt, bottom=8pt,
}

\newtcolorbox{cajaEjemplo}{
    enhanced, breakable,
    arc=6pt, outer arc=6pt,
    colback=cyanNeon!5!white,
    colframe=cyanNeon!70!azulNoche, boxrule=1pt,
    borderline west={4pt}{0pt}{cyanNeon},
    fonttitle=\bfseries\small\color{fondoOscuro},
    title={\faIcon{code}~Ejemplo},
    attach boxed title to top left={yshift=-3mm, xshift=6mm},
    boxed title style={colback=cyanNeon, arc=4pt, boxrule=0pt},
    drop fuzzy shadow=azulNoche!20!white,
    left=10pt, right=10pt, top=8pt, bottom=8pt,
}

\newtcolorbox{cajaEjercicio}{
    enhanced, breakable,
    arc=6pt, outer arc=6pt,
    colback=azulMedio!7!white,
    colframe=azulMedio, boxrule=1pt,
    borderline west={4pt}{0pt}{azulMedio},
    fonttitle=\bfseries\small\color{white},
    title={\faIcon{pencil-alt}~Ejercicio},
    attach boxed title to top left={yshift=-3mm, xshift=6mm},
    boxed title style={colback=azulMedio, arc=4pt, boxrule=0pt},
    drop fuzzy shadow=azulNoche!20!white,
    left=10pt, right=10pt, top=8pt, bottom=8pt,
}

\newtcolorbox{cajaReto}{
    enhanced, breakable,
    arc=6pt, outer arc=6pt,
    colback=naranjaVivo!6!white,
    colframe=naranjaVivo, boxrule=1pt,
    borderline west={4pt}{0pt}{naranjaVivo},
    fonttitle=\bfseries\small\color{white},
    title={\faIcon{fire}~Reto},
    attach boxed title to top left={yshift=-3mm, xshift=6mm},
    boxed title style={colback=naranjaVivo, arc=4pt, boxrule=0pt},
    drop fuzzy shadow=naranjaVivo!20!white,
    left=10pt, right=10pt, top=8pt, bottom=8pt,
}

% ── Indicadores de dificultad ─────────────────────────────────────────────────
\newcommand{\facil}{\textcolor{verdeSignal}{\faIcon{circle}\,\textbf{Fácil}}}
\newcommand{\intermedio}{\textcolor{naranjaVivo}{\faIcon{circle}\,\textbf{Intermedio}}}
\newcommand{\dificil}{\textcolor{rojoAlerta}{\faIcon{circle}\,\textbf{Difícil}}}

% ─────────────────────────────────────────────────────────────────────────────
\renewcommand{\iconoBanda}{flask}
\fancyhead[L]{\footnotesize\color{grisTexto}\faIcon{flask}~LABORATORIO\_I\_FISICA}
\fancyhead[R]{\footnotesize\color{grisTexto}Sílabo}
\hypersetup{pdftitle={Laboratorio I de Física},pdfauthor={Prof. César Rodríguez}
  pdfsubject={LABORATORIO\_I\_FISICA}}

\begin{document}
\bandaTitulo[fondoOscuro]{Laboratorio I de Física}{Medición, análisis gráfico y propiedades de la materia: del error experimental al informe reproducible}

\begin{center}\small\color{grisTexto}Prof. César Rodríguez \textbar{} Universidad de Oriente, Núcleo de Sucre, Departamento de Física \textbar{} 6 de noviembre de 2026\end{center}

\section{\iconotexto{id-card}{Identificación}}

\begin{tabularx}{\linewidth}{@{}>{\raggedright\arraybackslash}p{4.6cm}X@{}}
\toprule
\textbf{Código} & LABORATORIO\_I\_FISICA \\
\textbf{Duración} & 10 semanas \\
\textbf{Grupos} & 2 grupos de trabajo \\
\textbf{Modalidad} & Laboratorio presencial en mesas de dos grupos con roles rotativos. Las cuatro primeras semanas son talleres de teoría de errores y análisis gráfico; a partir de la semana 5 cada sesión es una práctica experimental con entrega semanal por repositorio privado. \\
\textbf{Entrega} & GitHub por grupo, revisado al finalizar cada semana \\
\textbf{Dedicación} & 3 horas de laboratorio + 2 horas de análisis e informe \\
\textbf{Versión} & 1 \\
\bottomrule
\end{tabularx}

\section{\iconotexto{bullseye}{Objetivo general}}

Formar al estudiante en el ciclo completo de una medición: elegir el instrumento, tomar datos crudos trazables, expresar el resultado con su incertidumbre, analizarlo gráficamente y comunicarlo en un informe que otro pueda reproducir. El curso usa las propiedades de la materia, la elasticidad, las oscilaciones y el enfriamiento como casos concretos para ejercitar ese ciclo, no como una colección de recetas aisladas.

\section{\iconotexto{tasks}{Resultados de aprendizaje}}

\begin{enumerate}[leftmargin=1.4em,itemsep=2pt,topsep=3pt]
  \item RA1: Expresar una medición con su incertidumbre, su unidad y sus cifras significativas, distinguiendo apreciación, exactitud y precisión.
  \item RA2: Propagar incertidumbres en magnitudes directas e indirectas y decidir cuándo una discrepancia es significativa.
  \item RA3: Graficar, linealizar y ajustar por mínimos cuadrados, reportando pendiente, intercepto, sus incertidumbres y los residuales.
  \item RA4: Determinar experimentalmente propiedades de la materia (volumen, densidad y centro de masa) y contrastarlas con un método independiente.
  \item RA5: Verificar leyes físicas extrayendo sus parámetros de un ajuste lineal, no del promedio de cocientes punto a punto.
  \item RA6: Redactar un informe de laboratorio reproducible, con datos crudos, script de análisis y figuras trazables a una tabla.
  \item RA7: Trabajar en equipo con roles rotativos, cumpliendo las normas de seguridad y el cuidado del instrumental.
\end{enumerate}

\section{\iconotexto{sitemap}{Estructura por fases}}

\begin{tcolorbox}[cajaContenido,title={Fase 1 --- Medición, error y análisis gráfico}]Semanas 1 a 4. Aprender a medir antes de medir cosas difíciles: apreciación, tipos de error, propagación, gráficas, linealización y ajuste por mínimos cuadrados, ejercitados con datos propios.
\end{tcolorbox}

\begin{tcolorbox}[cajaContenido,title={Fase 2 --- Propiedades de la materia y equilibrio}]Semanas 5 a 7. Medir volumen, densidad y centro de masa con dos métodos independientes, y usar la discrepancia entre ambos como criterio de control del montaje.
\end{tcolorbox}

\begin{tcolorbox}[cajaContenido,title={Fase 3 --- Elasticidad, oscilaciones y calor}]Semanas 8 a 10. Extraer parámetros físicos de ajustes lineales: constante elástica, aceleración de la gravedad y constante de enfriamiento, con sus incertidumbres y sus modelos de linealización.
\end{tcolorbox}

\section{\iconotexto{clipboard-check}{Sistema de evaluación}}

\begin{tabularx}{\linewidth}{@{}X r@{}}
\toprule
Concepto & Peso (\%)\\
\midrule
Talleres y prácticas (semanas 1 a 10) & 60,0 \\
Controles de fase (semanas 4 y 7) & 20,0 \\
Evaluación final práctica y portafolio (semana 10) & 20,0 \\
\midrule
\textbf{Total} & \textbf{100,0}\\
\bottomrule
\end{tabularx}

\begin{cajaRecuerda}La escala de calificación es de 10 puntos y la nota mínima aprobatoria es 5. La nota del curso se calcula ponderando los conceptos de la tabla; no se promedian notas de distinta escala.
\end{cajaRecuerda}

\section{\iconotexto{code-branch}{Entrega y revisión semanal}}

\begin{itemize}[leftmargin=1.4em,itemsep=2pt,topsep=3pt]
  \item Plataforma: GitHub.
  \item Modalidad: Un repositorio privado por grupo, con el docente como colaborador de lectura. El repositorio es la entrega; no se reciben informes por correo ni impresos, salvo indicación expresa del docente.
  \item Cadencia: El docente revisa el repositorio al finalizar cada semana y devuelve observaciones antes de la semana siguiente.
  \item Corte: Domingo a las 23:59 de la semana correspondiente. Lo entregado después del corte se corrige con la penalización prevista en las políticas.
\end{itemize}

\textbf{Estructura del repositorio}

\begin{itemize}[leftmargin=1.4em,itemsep=2pt,topsep=3pt]
  \item README.md con integrantes, roles de la fase y tabla de semanas entregadas.
  \item semana\_NN/informe.pdf con el informe de la práctica o el taller resuelto.
  \item semana\_NN/datos/ con los CSV crudos, uno por ensayo, sin editar.
  \item semana\_NN/analisis/ con los scripts que reproducen cada figura.
  \item semana\_NN/figuras/ con las gráficas en PDF o PNG generadas por los scripts.
  \item semana\_NN/README.md con lo que se hizo, lo que no se pudo verificar y por qué.
\end{itemize}

\textbf{Reglas de entrega}

\begin{itemize}[leftmargin=1.4em,itemsep=2pt,topsep=3pt]
  \item Los datos crudos no se editan nunca; toda corrección vive en el script de análisis.
  \item Cada cifra del informe debe salir de una tabla de datos o de un script entregado.
  \item Los scripts se ejecutan de punta a punta sin pasos manuales intermedios.
  \item Ninguna tabla se entrega sin unidades ni sin la resolución del instrumento.
  \item El informe es grupal, pero las conclusiones se firman individualmente.
  \item Un resultado sin incertidumbre no se acepta como resultado.
\end{itemize}

\begin{cajaRecuerda}La entrega se considera válida solo cuando está en el repositorio del grupo antes del corte. Un enlace, una captura suelta o un archivo enviado por mensajería no sustituyen el repositorio.\end{cajaRecuerda}

\section{\iconotexto{server}{Laboratorio}}

Mesas con regla milimetrada, calibrador, tornillo micrométrico, balanza, probeta, cronómetro, termómetro digital, soportes, resortes, masas calibradas y cuerpos sólidos de distintos materiales y tamaños.

Laboratorio físico de mesa. Cada grupo dispone de un juego de instrumentos y de una estación de análisis con Python para procesar sus propios datos.

\begin{itemize}[leftmargin=1.4em,itemsep=2pt,topsep=3pt]
  \item No se energiza ningún equipo sin revisión previa del docente.
  \item Los líquidos se manipulan lejos de las balanzas y de las conexiones eléctricas.
  \item Cada instrumento se devuelve limpio, seco y en su estuche al terminar la sesión.
  \item El material roto se reporta de inmediato; ocultarlo agrava la sanción.
  \item Ningún dato se inventa ni se completa de memoria después de la sesión.
\end{itemize}

\section{\iconotexto{balance-scale}{Políticas de uso}}

\begin{itemize}[leftmargin=1.4em,itemsep=2pt,topsep=3pt]
  \item Seguridad: no se energiza ningún equipo ni se manipula vidrio caliente sin autorización y sin el equipo de protección indicado.
  \item Honestidad académica: los datos se toman en la sesión y se declaran tal como se obtuvieron; fabricar o alterar datos anula la entrega y se reporta.
  \item Puntualidad: la entrega cierra el domingo a las 23:59; lo entregado después se corrige con la penalización prevista.
  \item Cuidado del material: cada instrumento se devuelve limpio y en su estuche; el material roto se reporta de inmediato.
  \item Trabajo en equipo: todos los integrantes firman las conclusiones y deben poder explicar cualquier parte del informe.
\end{itemize}

\section{\iconotexto{book}{Bibliografía}}

\begin{itemize}[leftmargin=1.4em,itemsep=2pt,topsep=3pt]
  \item Baird, D. C.. Experimentación: una introducción a la teoría de mediciones y al diseño de experimentos. Prentice Hall, 1991.
  \item Taylor, J. R.. An Introduction to Error Analysis. University Science Books, 1997.
  \item Sears, F. y Zemansky, M.. Física Universitaria. Pearson Educación, edición reciente.
  \item Serway, R. y Jewett, J.. Física para ciencias e ingeniería. Cengage Learning, edición reciente.
  \item Departamento de Física, Universidad de Oriente. Guía de laboratorio de Física I. Material del curso, 2026.
\end{itemize}
\end{document}
